% !PRINTED 429
\weberpage{429}
%<<CONT
so folgt hieraus die Gleichung 6\textsuperscript{ten} Grades für $U$:
\weq{7}{U^6 - 10\,z\,U^3 + z^2\,U + 5\,z^2 = 0}
als \gesperrt{Resolvente} 6\textsuperscript{ten} \gesperrt{Grades} der Ikosaëdergleichung.

Setzen wir nach (3) und (4)
\weq{8}{\begin{aligned}
\sqrt{w_\infty} &= \sqrt{5}\,x_1\,x_2 \\
\sqrt{w_r} &= \varepsilon^r x_1^2 + x_1 x_2 - \varepsilon^{-r} x_2^2,
\end{aligned}}
so ergeben sich daraus die folgenden Relationen:
\weq{9}{\begin{aligned}
\sqrt{5\,w_\infty} &= \sqrt{w_0} + \sqrt{w_1} + \sqrt{w_2} + \sqrt{w_3} + \sqrt{w_4} \\
0 &= \sqrt{w_0} + \varepsilon^2\sqrt{w_1} + \varepsilon^4\sqrt{w_2} + \varepsilon\sqrt{w_3} + \varepsilon^3\sqrt{w_4} \\
0 &= \sqrt{w_0} + \varepsilon^3\sqrt{w_1} + \varepsilon\sqrt{w_2} + \varepsilon^4\sqrt{w_3} + \varepsilon^2\sqrt{w_4}
\end{aligned}}
\weq{10}{\frac{x_1\sqrt{5}}{x_2} = \frac{\sqrt{w_0} + \varepsilon^4\sqrt{w_1} + \varepsilon^3\sqrt{w_2} + \varepsilon^2\sqrt{w_3} + \varepsilon\sqrt{w_4}}{\sqrt{w_\infty}},}
und in den Gleichungen (9), (10) können, da sie homogen sind, an Stelle der $w$ auch die $U$ gesetzt werden.

Auf die Gleichungen 6\textsuperscript{ten} Grades, deren Wurzeln den Relationen (9) genügen, hat zuerst \gesperrt{Jacobi} hingewiesen; man nennt sie daher \gesperrt{Jacobi}'sche Gleichungen 6\textsuperscript{ten} Grades\footnote{Suite des notices sur les fonctions elliptiques. \gesperrt{Crelle's} Journal, Bd. III (1828); Werke, Bd. I, S. 261.}.

\begin{center}
\S.\,114.\\[4pt]
\gesperrt{Ueber die Lösung der Ikosaëdergleichung durch transcendente Functionen.}
\end{center}

Wir wollen hier, ohne auf Beweise oder nähere Ausführungen einzugehen, kurz auf die Beziehungen hinweisen, die zwischen den durchgeführten Untersuchungen und der Theorie gewisser transcendenter Functionen, besonders der Theorie der elliptischen Functionen und der Theorie der hypergeometrischen Reihen bestehen, woraus sich die Auflösung der Ikosaëdergleichung und damit aller ihrer Resolventen durch transcendente Functionen ergiebt. Dies führt aus dem Gebiete der Algebra hinaus, steht aber zu ihr etwa in derselben Beziehung, wie wenn man bei der Berechnung von Wurzelgrössen Logarithmen, oder bei den cubischen Gleichungen die Winkeltheilung anwendet.
